\subsection{Framework Architecture}

The \texttt{bioimage\_embed} framework consists of four main components: data handling, model architectures, training infrastructure, and deployment tools (Figure~\ref{fig:architecture}). We describe each component below.

\subsection{Data Handling}

\subsubsection{File Format Support}

We implement custom PyTorch \texttt{Dataset} classes for common microscopy formats:

\begin{itemize}
    \item \textbf{TiffDataset:} Handles multi-channel TIFF stacks with lazy loading
    \item \textbf{NgffDataset:} Supports OME-NGFF (Next Generation File Format) with zarr backends
    \item \textbf{DatasetGlob:} Flexible file-pattern matching for heterogeneous datasets
\end{itemize}

All datasets inherit from a common base class that provides:
\begin{itemize}
    \item Channel normalization (per-channel z-scoring or percentile-based)
    \item Configurable augmentations via Albumentations
    \item Metadata extraction and preservation
\end{itemize}

\subsubsection{Integration with Public Repositories}

We provide pre-configured loaders for:
\begin{itemize}
    \item \textbf{BioImage Data Resource (IDR):} Direct access to curated microscopy datasets
    \item \textbf{Cell Image Library:} Standardized cell and tissue images
    \item \textbf{Cell Painting:} High-throughput phenotypic screening data
\end{itemize}

\subsection{Model Architectures}

\subsubsection{Variational Autoencoders (VAE)}

Our VAE implementation follows the standard formulation with encoder $q_\phi(z|x)$ and decoder $p_\theta(x|z)$:

\begin{equation}
\mathcal{L}_{\text{VAE}} = \mathbb{E}_{q_\phi(z|x)}[\log p_\theta(x|z)] - \beta \cdot D_{KL}(q_\phi(z|x) \| p(z))
\end{equation}

where $\beta$ controls the regularization strength ($\beta$-VAE variant).

\textbf{Encoder backbones:}
\begin{itemize}
    \item ResNet18, ResNet50 (from torchvision)
    \item Vision Transformer (ViT) with patch embeddings
\end{itemize}

\textbf{Decoder architectures:}
\begin{itemize}
    \item Transposed convolution decoder (symmetric to ResNet encoder)
    \item ViT-based decoder with patch reconstruction
\end{itemize}

\subsubsection{Vector Quantized VAE (VQ-VAE)}

We implement VQ-VAE with learnable codebook $\mathcal{C} = \{c_k\}_{k=1}^K$:

\begin{equation}
z_q = \argmin_{c_k \in \mathcal{C}} \|z_e - c_k\|_2
\end{equation}

where $z_e$ is the encoder output and $z_q$ is the quantized latent code.

\subsubsection{Masked Autoencoders (MAE)}

Following \citet{he2022masked}, we implement masked autoencoders with:
\begin{itemize}
    \item Random patch masking (75\% masking ratio)
    \item ViT encoder operating only on visible patches
    \item Lightweight decoder reconstructing masked patches
\end{itemize}

\subsubsection{Supervised Learning Mode}

For datasets with labels, we support supervised and semi-supervised learning:

\begin{equation}
\mathcal{L}_{\text{total}} = \mathcal{L}_{\text{reconstruction}} + \lambda \cdot \mathcal{L}_{\text{classification}}
\end{equation}

This enables joint optimization of reconstruction quality and downstream task performance.

\subsection{Training Infrastructure}

\subsubsection{Configuration Management}

We use Hydra for hierarchical configuration with Pydantic dataclasses for type safety:

\begin{verbatim}
@dataclass
class ModelConfig:
    architecture: str  # "vae", "vqvae", "mae"
    latent_dim: int
    beta: float = 1.0
    encoder_backbone: str = "resnet18"
\end{verbatim}

All hyperparameters are validated at runtime and saved with each checkpoint.

\subsubsection{Reproducibility}

Each training run saves:
\begin{itemize}
    \item Full resolved configuration (YAML)
    \item Git commit SHA
    \item Dataset version/hash
    \item Random seeds (PyTorch, NumPy, Python)
    \item Environment snapshot (conda/pip freeze)
\end{itemize}

\subsubsection{Logging and Monitoring}

Integrated support for:
\begin{itemize}
    \item TensorBoard (local logging)
    \item Weights \& Biases (cloud logging)
    \item Custom callbacks for reconstruction visualization
\end{itemize}

\subsection{Deployment}

\subsubsection{Model Export}

Trained models can be exported to:
\begin{itemize}
    \item ONNX format for deployment in non-Python environments
    \item HuggingFace Hub for sharing and versioning
    \item PyTorch checkpoint with full training state
\end{itemize}

\subsubsection{Command-Line Interface}

The CLI provides high-level commands:

\begin{verbatim}
# Train a VAE on Cell Painting data
bioimage_embed train --data cellpainting \
    --model vae --latent-dim 128

# Encode new images
bioimage_embed encode --checkpoint model.ckpt \
    --input images/ --output embeddings.npy
\end{verbatim}

\subsection{Software Engineering}

The codebase follows modern Python best practices:
\begin{itemize}
    \item Type hints throughout (mypy strict mode)
    \item Pre-commit hooks (black, isort, flake8)
    \item Unit tests with pytest
    \item Continuous integration via GitHub Actions
    \item Semantic versioning
\end{itemize}
